\documentclass[a4paper,12pt]{article}

\usepackage[T2A]{fontenc}
\usepackage[utf8]{inputenc}
\usepackage[russian]{babel}
\usepackage{amsmath,amssymb,amsthm,mathtools}
\usepackage{helvet}
\renewcommand{\familydefault}{\sfdefault}
\usepackage{icomma}
\usepackage{geometry}
\usepackage{microtype}
\usepackage{tikz}
\usetikzlibrary{calc,arrows.meta,positioning}
\usepackage{pgfplots}
\pgfplotsset{compat=1.18}
\usepackage{tabularx,booktabs,longtable}
\usepackage{enumitem}
\usepackage{hyperref}
\usepackage{parskip}
\usepackage{fancyhdr}
\usepackage{setspace}
\usepackage{xcolor}

\geometry{margin=2.1cm}
\onehalfspacing

\definecolor{accent}{HTML}{008080}
\definecolor{darkaccent}{HTML}{004D4D}
\definecolor{textgray}{HTML}{333333}
\definecolor{softgray}{HTML}{737373}
\definecolor{lightaccent}{HTML}{DDEEEE}

\color{textgray}

\hypersetup{
    colorlinks=true,
    linkcolor=darkaccent,
    urlcolor=darkaccent
}

\setlist[itemize]{
    leftmargin=1.5em,
    itemsep=0.25em,
    topsep=0.35em
}

\setlist[enumerate]{
    leftmargin=1.7em,
    itemsep=0.55em,
    topsep=0.45em
}

\pagestyle{fancy}
\fancyhf{}
\fancyhead[L]{\small\textcolor{softgray}{Прикладные задачи ЕГЭ}}
\fancyhead[R]{\small\textcolor{softgray}{Тимофей Васильченко}}
\fancyfoot[C]{\small\textcolor{softgray}{\thepage}}
\renewcommand{\headrulewidth}{0.3pt}
\renewcommand{\headrule}{\hbox to\headwidth{\color{lightaccent}\leaders\hrule height \headrulewidth\hfill}}

\newcommand{\sectionline}{%
    \par\vspace{0.15em}
    {\color{lightaccent}\hrule height 0.7pt}
    \vspace{0.7em}
}

\newcommand{\key}[1]{\textcolor{darkaccent}{\textbf{#1}}}
\newcommand{\dd}{\,\mathrm{d}}

\begin{document}

% ------------------------------------------------------------
% ТИТУЛЬНЫЙ ЛИСТ
% ------------------------------------------------------------

\begin{titlepage}
\thispagestyle{empty}

\begin{center}

\vspace*{3.0cm}

{\Large\textcolor{softgray}{Чек-лист}\par}

\vspace{0.7cm}

{\fontsize{26}{32}\selectfont
\textbf{\textcolor{darkaccent}{Прикладные задачи}}\par}

\vspace{0.25cm}

{\Large
формулы, размерности, степени\\
и математические модели\par}

\vspace{1.8cm}

{\large
Лёвин Артём Александрович\\[0.25cm]
эксклюзивно для\\[0.25cm]
\textbf{Тимофея Васильченко}\par}

\vspace{1.3cm}

{\large \today}

\end{center}

\vfill

\begin{center}
\begin{tikzpicture}[x=1cm,y=1cm]
    \draw[
        darkaccent,
        line width=1.1pt
    ]
    (0,0)
    .. controls (3.1,0.55) and (6.1,-0.45) ..
    (9.0,0)
    .. controls (11.2,0.34) and (13.2,0.20) ..
    (15.0,-0.05);
\end{tikzpicture}
\end{center}

\vspace{0.8cm}

\end{titlepage}

% ------------------------------------------------------------
% ОСНОВНОЙ ЧЕК-ЛИСТ
% ------------------------------------------------------------

\section*{\textcolor{darkaccent}{1. Что проверяет прикладная задача}}

\sectionline

Прикладная задача обычно содержит формулу из физики, техники, экономики
или другого практического сюжета. Главная трудность состоит не в самой
формуле, а в последовательном переводе условия на математический язык.

\key{Перед вычислениями проверьте четыре пункта:}

\begin{enumerate}
    \item Что именно требуется найти?
    \item Какие величины известны?
    \item Все ли величины записаны в совместимых единицах измерения?
    \item Можно ли сначала выразить искомую величину буквенно и только затем
    подставлять числа?
\end{enumerate}

\begin{center}
\begin{tikzpicture}[
    node distance=0.8cm and 0.65cm,
    every node/.style={font=\small},
    box/.style={
        draw=lightaccent,
        rounded corners=2pt,
        align=center,
        inner sep=6pt,
        text width=3.1cm
    },
    arr/.style={-{Latex[length=2mm]},darkaccent,line width=0.8pt}
]
\node[box] (a) {\textbf{Текст}\\условия};
\node[box,right=of a] (b) {\textbf{Модель}\\и формула};
\node[box,right=of b] (c) {\textbf{Преобразование}\\в символах};
\node[box,right=of c] (d) {\textbf{Расчёт}\\и отбор ответа};

\draw[arr] (a) -- (b);
\draw[arr] (b) -- (c);
\draw[arr] (c) -- (d);
\end{tikzpicture}
\end{center}

\section*{\textcolor{darkaccent}{2. Размерности: сначала привести к общим единицам}}

\sectionline

Если в одной формуле используются длины, массы, время или другие величины,
они должны быть выражены в согласованных единицах.

Например,
\[
1\text{ мм}=10^{-3}\text{ м},
\qquad
3\text{ мм}=3\cdot10^{-3}\text{ м}.
\]

\key{Типичная ошибка:} подставить в одну формулу \(10\) метров и \(3\)
миллиметра как числа \(10\) и \(3\). Такой расчёт теряет физический смысл.

\subsection*{\textcolor{accent}{Пример: тепловое расширение рельса}}

Длина рельса описывается формулой
\[
L=L_0(1+\alpha T),
\]
где \(L_0\) --- первоначальная длина, \(\alpha\) --- коэффициент линейного
расширения, \(T\) --- изменение температуры.

Пусть
\[
L_0=10\text{ м},
\qquad
\alpha=1,2\cdot10^{-5},
\qquad
\Delta L=3\text{ мм}.
\]

Сначала переводим удлинение:
\[
\Delta L=3\cdot10^{-3}\text{ м}.
\]

По смыслу удлинения
\[
\Delta L=L-L_0,
\]
следовательно,
\[
L=L_0+\Delta L.
\]

Теперь одна и та же конечная длина \(L\) выражена двумя способами:
\[
L_0(1+\alpha T)=L_0+\Delta L.
\]

Раскрываем скобки:
\[
L_0+L_0\alpha T=L_0+\Delta L.
\]

Отсюда
\[
L_0\alpha T=\Delta L,
\]
поэтому
\[
\boxed{T=\frac{\Delta L}{L_0\alpha}}.
\]

Для удобства вычислений преобразуем
\[
1,2\cdot10^{-5}=12\cdot10^{-6}.
\]

Тогда
\[
T=
\frac{3\cdot10^{-3}}
{10\cdot12\cdot10^{-6}}
=25.
\]

\[
\boxed{T=25}
\]

\key{Главная идея:} сначала вывести удобную рабочую формулу, затем заниматься
числами.

\section*{\textcolor{darkaccent}{3. Что означает символ \(\Delta\)}}

\sectionline

Символ \(\Delta\) обозначает \textbf{изменение величины}.

Для длины:
\[
\boxed{\Delta L=L_{\text{нов}}-L_{\text{стар}}}.
\]

Если старая длина равна \(L_0\), а новая --- \(L\), то
\[
\Delta L=L-L_0.
\]

Отсюда
\[
L=L_0+\Delta L.
\]

\key{Проверка здравого смысла:} при положительном удлинении новая длина
обязана быть больше первоначальной.

Поэтому выражение
\[
L=L_0-\Delta L
\]
при положительном удлинении противоречило бы смыслу задачи.

\section*{\textcolor{darkaccent}{4. Как правильно выражать переменную из формулы}}

\sectionline

Пусть
\[
A=B+C\alpha T.
\]

Если требуется найти \(T\), действуйте последовательно.

Сначала устраняем слагаемые:
\[
A-B=C\alpha T.
\]

Затем устраняем множители:
\[
T=\frac{A-B}{C\alpha}.
\]

\key{Важно:} перенос через знак равенства --- это сокращённое описание
одинаковой операции над обеими частями уравнения.

Если величина \(\alpha\) \textbf{умножается} на \(T\), то при выражении \(T\)
она окажется в знаменателе:
\[
\alpha T=M
\quad\Longrightarrow\quad
T=\frac{M}{\alpha},
\qquad \alpha\ne0.
\]

\subsection*{\textcolor{accent}{Функциональная запись}}

Запись
\[
L(T)
\]
означает: \textit{величина \(L\) зависит от \(T\)}.

Скобки в такой записи не являются дополнительным множителем.
При обычном алгебраическом преобразовании можно обозначить значение функции
просто буквой \(L\), если это не создаёт неоднозначности.

% ------------------------------------------------------------

\newpage

\section*{\textcolor{darkaccent}{5. Научная запись чисел}}

\sectionline

Для больших и малых чисел удобно использовать вид
\[
a\cdot10^n,
\qquad
1\le |a|<10.
\]

Однако при ручных расчётах бывает полезно временно передвинуть запятую так,
чтобы коэффициент стал целым.

Например,
\[
1,2\cdot10^{-5}=12\cdot10^{-6},
\]
\[
5,7\cdot10^{-8}=57\cdot10^{-9},
\]
\[
4,5\cdot10^{-3}=45\cdot10^{-4}.
\]

Правило:
\[
a\cdot10^n=(10a)\cdot10^{n-1}.
\]

Если запятая смещается вправо на один разряд, показатель степени десятки
уменьшается на \(1\).

\section*{\textcolor{darkaccent}{6. Степени десятки при умножении и делении}}

\sectionline

Основные свойства:
\[
10^a\cdot10^b=10^{a+b},
\]
\[
\frac{10^a}{10^b}=10^{a-b},
\]
\[
(10^a)^b=10^{ab}.
\]

Особенно внимательно работайте с отрицательными показателями:
\[
\frac{10^{-3}}{10^{-6}}
=
10^{-3-(-6)}
=
10^3.
\]

\key{Типичная ошибка:} потерять второй минус в выражении
\[
a-(-b)=a+b.
\]

\section*{\textcolor{darkaccent}{7. Большие расчёты: разложение на простые множители}}

\sectionline

Если формула содержит большие коэффициенты, степени и корни, обычный
десятичный расчёт часто становится громоздким.

В таких ситуациях числа полезно раскладывать на \textbf{простые множители}.

Например,
\[
912=2^4\cdot3\cdot19,
\]
\[
57=3\cdot19,
\]
\[
16=2^4.
\]

Тогда выражение
\[
\frac{912\cdot16}{57}
\]
превращается в
\[
\frac{
2^4\cdot3\cdot19\cdot2^4
}{
3\cdot19
}
=
2^8.
\]

Такой вид особенно удобен, если следующим шагом потребуется извлекать корень.

\key{Простое число} имеет ровно два положительных делителя: \(1\) и само
себя.

Например, \(2,3,5,7,11,19\) --- простые числа.

Числа \(4,10,12,16\) --- составные, поэтому при полном разложении их следует
раскладывать дальше.

\section*{\textcolor{darkaccent}{8. Извлечение корня через степени}}

\sectionline

Если
\[
T^4=A,
\qquad
A\ge0,
\]
и по смыслу задачи \(T\ge0\), то
\[
T=A^{1/4}.
\]

Например,
\[
T^4=2^8\cdot10^{12}.
\]

Возводим обе части в степень \(\frac14\):
\[
T=
\left(
2^8\cdot10^{12}
\right)^{1/4}.
\]

Используем свойства степеней:
\[
T=
2^{8/4}\cdot10^{12/4}
=
2^2\cdot10^3.
\]

Следовательно,
\[
\boxed{T=4000}.
\]

\key{Почему такой способ удобен:} показатели степеней делятся на показатель
корня.

В общем виде:
\[
\sqrt[n]{a^m}=a^{m/n},
\]
если выражение определено в рассматриваемой области.

Для чётного \(n\) арифметический корень всегда неотрицателен.

\section*{\textcolor{darkaccent}{9. Пример большой прикладной формулы}}

\sectionline

Рассмотрим закон вида
\[
P=\sigma S T^4.
\]

Пусть
\[
P=9,12\cdot10^{25},
\qquad
\sigma=5,7\cdot10^{-8},
\qquad
S=\frac{10^{20}}{16}.
\]

Требуется найти \(T\).

Сначала выражаем четвёртую степень температуры:
\[
T^4=\frac{P}{\sigma S}.
\]

Преобразуем десятичные коэффициенты:
\[
P=912\cdot10^{23},
\qquad
\sigma=57\cdot10^{-9}.
\]

Тогда
\[
T^4=
\frac{
912\cdot10^{23}
}{
57\cdot10^{-9}\cdot\dfrac{10^{20}}{16}
}.
\]

Деление на дробь заменяем умножением:
\[
T^4=
\frac{
912\cdot16\cdot10^{23}
}{
57\cdot10^{-9}\cdot10^{20}
}.
\]

Разложим коэффициенты:
\[
912=2^4\cdot3\cdot19,
\qquad
16=2^4,
\qquad
57=3\cdot19.
\]

Поэтому
\[
T^4
=
2^8\cdot10^{23-(-9)-20}
=
2^8\cdot10^{12}.
\]

Следовательно,
\[
T=
\left(2^8\cdot10^{12}\right)^{1/4}
=
2^2\cdot10^3
=
\boxed{4000}.
\]

\key{Алгоритм большого расчёта:}

\begin{enumerate}
    \item Выразить искомую величину.
    \item Убрать неудобные десятичные коэффициенты.
    \item Разложить целые коэффициенты на простые множители.
    \item Отдельно обработать степени десятки.
    \item Сократить одинаковые множители.
    \item Только после этого извлекать корень или выполнять финальное вычисление.
\end{enumerate}

% ------------------------------------------------------------

\newpage

\section*{\textcolor{darkaccent}{10. Текстовая задача с траекторией}}

\sectionline

В некоторых прикладных задачах формула задаёт высоту объекта в зависимости
от расстояния.

Например,
\[
y=-0,01x^2+x.
\]

Пусть высота стены равна \(8\) м, а камень должен пройти минимум на \(1\) м
выше стены.

Требуемая высота:
\[
y=8+1=9.
\]

Подставляем в формулу:
\[
-0,01x^2+x=9.
\]

Умножаем уравнение на \(100\):
\[
-x^2+100x=900,
\]
\[
x^2-100x+900=0.
\]

Корни:
\[
x_1=10,
\qquad
x_2=90.
\]

Если требуется \textbf{наибольшее} расстояние, выбираем
\[
\boxed{90}.
\]

\begin{center}
\begin{tikzpicture}
\begin{axis}[
    width=0.82\linewidth,
    height=7.0cm,
    axis lines=middle,
    unit vector ratio=1 1,
    xmin=0,
    xmax=100,
    ymin=0,
    ymax=28,
    samples=250,
    xlabel={$x$},
    ylabel={$y$},
    xtick={0,10,50,90,100},
    ytick={0,9,25},
    clip=false,
    tick label style={font=\small},
    label style={font=\small},
    every axis plot/.append style={line width=1.1pt}
]
\addplot[darkaccent,domain=0:100] {-0.01*x^2+x};
\addplot[softgray,dashed,domain=0:100] {9};

\addplot[
    only marks,
    mark=*,
    mark size=2pt
]
coordinates {(10,9) (90,9)};

\node[anchor=south west,font=\small] at (axis cs:52,25)
{$y=-0,01x^2+x$};

\node[anchor=south west,font=\small] at (axis cs:62,9.5)
{$y=9$};
\end{axis}
\end{tikzpicture}
\end{center}

\key{Важно:} оба корня могут быть математически допустимы, однако вопрос
задачи определяет, какой из них нужно записать в ответ.

\section*{\textcolor{darkaccent}{11. Слова «не менее», «не более», «наибольшее»}}

\sectionline

В прикладной задаче слова условия определяют математическое ограничение:

\[
\text{«не менее }a\text{»}
\quad\Longleftrightarrow\quad
x\ge a,
\]
\[
\text{«не более }a\text{»}
\quad\Longleftrightarrow\quad
x\le a.
\]

Однако иногда искомое граничное значение определяется равенством.

Например, если требуется найти максимальное расстояние, при котором высота
остаётся не менее \(9\) м, граница допустимого участка находится из
\[
y=9.
\]

После решения уравнения обязательно возвращайтесь к вопросу задачи и
проверяйте, какой корень соответствует требуемой границе.

\section*{\textcolor{darkaccent}{12. Финальный алгоритм решения}}

\sectionline

\begin{enumerate}
    \item \key{Прочитайте вопрос.} Сразу зафиксируйте искомую величину.
    \item \key{Выпишите данные.} Отдельно укажите формулу.
    \item \key{Приведите единицы.} Нельзя смешивать, например, метры и миллиметры.
    \item \key{Расшифруйте смысл величин.} Например,
    \[
    \Delta L=L-L_0.
    \]
    \item \key{Составьте математическую модель.}
    \item \key{Выразите искомую величину буквенно.}
    \item \key{Подготовьте числа к вычислению.} Уберите неудобные десятичные дроби.
    \item \key{При больших расчётах разложите коэффициенты на простые множители.}
    \item \key{Используйте свойства степеней.}
    \item \key{Проверьте ограничения и смысл ответа.}
\end{enumerate}

\section*{\textcolor{darkaccent}{13. Типичные ошибки}}

\sectionline

\begin{itemize}
    \item Смешивать различные единицы измерения.
    \item Неверно понимать изменение величины:
    \[
    \Delta L=L_{\text{нов}}-L_{\text{стар}}.
    \]
    \item Путать сложение и умножение при выражении переменной.
    \item Считать запись \(L(T)\) произведением \(L\cdot T\).
    \item Слишком рано подставлять числа и получать громоздкий расчёт.
    \item Терять знак при работе с отрицательными показателями степеней.
    \item Оставлять составные числа, когда требуется дальнейшая работа со степенями.
    \item Извлекать корень из произведения, не приведя выражение к удобному степенному виду.
    \item Записывать любой найденный корень без проверки вопроса задачи.
    \item Игнорировать физические ограничения: длина, расстояние, площадь и абсолютная температура
    в стандартных прикладных задачах не могут быть отрицательными.
\end{itemize}

% ------------------------------------------------------------
% ТРЕНИРОВОЧНЫЙ БЛОК
% ------------------------------------------------------------

\newpage
\thispagestyle{fancy}

\section*{\textcolor{darkaccent}{Тренировочный блок}}

\sectionline

\textbf{Формат: Путь А.}

Выполняйте задачи последовательно. В расчётных задачах старайтесь сначала
получить рабочую формулу в буквенном виде.

\subsection*{\textcolor{accent}{Средний уровень}}

\begin{enumerate}[label=\textbf{\arabic*.}]

\item
Длина металлического стержня при нагревании описывается формулой
\[
L=L_0(1+\alpha T).
\]
Известно, что
\[
L_0=8\text{ м},
\qquad
\alpha=1,25\cdot10^{-5},
\]
а стержень удлинился на \(2\) мм.

Найдите \(T\).

\item
Величина \(Q\) определяется формулой
\[
Q=ab^2.
\]
Найдите \(b>0\), если
\[
Q=1,8\cdot10^5,
\qquad
a=5\cdot10^{-3}.
\]

\item
Высота тела над землёй задаётся формулой
\[
y=-0,02x^2+2x.
\]
На каком \textbf{наибольшем} расстоянии \(x\) высота тела равна \(32\) м?

\end{enumerate}

\subsection*{\textcolor{accent}{Выше среднего}}

\begin{enumerate}[label=\textbf{\arabic*.},start=4]

\item
По формуле
\[
P=\sigma ST^4
\]
найдите \(T>0\), если
\[
P=8,1\cdot10^{17},
\qquad
\sigma=5\cdot10^{-7},
\qquad
S=2\cdot10^8.
\]

\item
Металлическая деталь имела длину \(12\) м и после нагревания удлинилась
на \(7,2\) мм. Коэффициент линейного расширения равен
\[
\alpha=2\cdot10^{-5}.
\]
Используя формулу
\[
L=L_0(1+\alpha T),
\]
найдите изменение температуры \(T\).

\item
Упростите выражение
\[
\frac{
4,8\cdot10^{-7}
}{
1,2\cdot10^{-10}
}.
\]

Ответ представьте в виде \(a\cdot10^n\), где \(1\le a<10\).

\item
Найдите положительное значение \(x\), если
\[
x^4=2^{12}\cdot5^8.
\]

\item
Высота снаряда описывается формулой
\[
h=-0,01x^2+1,2x.
\]
Высота стены равна \(20\) м. Для безопасного прохождения снаряд должен
находиться минимум на \(5\) м выше стены.

Найдите \textbf{наибольшее} расстояние \(x\), на котором выполняется
граничное условие.

\item
По формуле
\[
E=kS^2T^4
\]
найдите \(T>0\), если
\[
E=2^{16}\cdot3^4\cdot10^{12},
\qquad
k=3^4,
\qquad
S=2^4.
\]

\end{enumerate}

\subsection*{\textcolor{accent}{Сложная задача}}

\begin{enumerate}[label=\textbf{\arabic*.},start=10]

\item
Мощность источника описывается формулой
\[
P=\sigma ST^4.
\]

Известно:
\[
P=3,456\cdot10^{28},
\qquad
\sigma=6\cdot10^{-8},
\qquad
S=9\cdot10^{20}.
\]

Найдите \(T>0\).

При вычислениях сначала преобразуйте десятичные коэффициенты, затем
разложите целые числа на простые множители.

\end{enumerate}

% ------------------------------------------------------------
% ОТВЕТЫ
% ------------------------------------------------------------

\newpage
\thispagestyle{fancy}

\section*{\textcolor{darkaccent}{Ответы}}

\sectionline

\renewcommand{\arraystretch}{1.4}

\begin{table}[ht]
\centering
\caption{Ответы к тренировочному блоку}
\begin{tabularx}{\linewidth}{@{}cX@{}}
\toprule
\textbf{№} & \textbf{Ответ} \\
\midrule

1 &
\[
T=20
\]
\\

2 &
\[
b=6000
\]
\\

3 &
\[
x=80
\]
\\

4 &
\[
T=3000
\]
\\

5 &
\[
T=30
\]
\\

6 &
\[
4\cdot10^3
\]
\\

7 &
\[
x=2^3\cdot5^2=200
\]
\\

8 &
\[
x=60+10\sqrt{11}
\]
\\

9 &
\[
T=2^2\cdot10^3=4000
\]
\\

10 &
\[
T=2\sqrt[4]{10^6}
=20\sqrt{10}
\]
\\

\bottomrule
\end{tabularx}
\end{table}

\vspace{1cm}

\section*{\textcolor{darkaccent}{Контроль перед завершением задачи}}

\sectionline

Перед тем как записать ответ, проверьте:

\begin{itemize}
    \item все ли величины приведены к согласованным единицам;
    \item верно ли выражена искомая переменная;
    \item не потерян ли минус в показателе степени;
    \item полностью ли разложены составные коэффициенты, если это было нужно;
    \item правильно ли извлечён корень;
    \item соответствует ли выбранный корень словам
    «наибольшее», «наименьшее», «не менее», «не более»;
    \item имеет ли ответ смысл в контексте задачи.
\end{itemize}

\end{document}